% from ega4-18-fr.tex
\oldpage[IV-4]{219}

Soit $\sh J$ l'Idéal cohérent de $\sh O_X$ définissant $Y$; l'hypothèse sur $\sh F$ signifie que $\operatorname{gr}_{\sh J}^*(\sh F)$ est un $\sh O_Y$-Module plat. On peut donc, pour tout $x\in Y$, appliquer (19.7.3) en prenant $A=\sh O_{X,x}$, $\mathfrak R=\sh J_x$, $M=\sh F_x$; l'hypothèse $a)$ étant remplie par hypothèse, la série de Poincaré $P_x(\sh F)(T)$ est donc égale à $R_x(T)(1-T)^{-n}$, où $R_x(T)$ est la série de Poincaré de $\operatorname{gr}_{\sh J_x}^*(\sh F_x)\otimes_{\sh O_{Y,x}}k(x)$. Or, chacun des $\sh O_Y$-Modules $\sh J^h\sh F/\sh J^{h+1}\sh F$ est plat et cohérent par hypothèse, donc localement libre (2.1.12), et puisque $Y$ est supposé connexe, chacun des $\sh J^h\sh F/\sh J^{h+1}\sh F$ est de rang constant dans $Y$, ce qui prouve que la série de Poincaré $R_x(T)$ est indépendante de $x\in Y$.
\end{corollary}

\begin{corollary}[19.7.5]
\label{IV.19.7.5}
Soient $X$ un préschéma localement noethérien, $Y$ un sous-préschéma fermé de $X$, $Z$ un sous-préschéma fermé de $Y$; on suppose $Y$ et $Z$ réguliers. Soit $\sh F$ un $\sh O_X$-Module cohérent.
\begin{enumerate}
  \item[(i)] Si $\sh F$ est normalement plat le long de $Y$ aux points de $Z$ (cf. (11.3.4) pour la définition de la platitude normale en un point), alors $\sh F$ est normalement plat le long de $Z$.
  \item[(ii)] Supposons en outre que $Z$ soit connexe et que $\sh O_X$ soit normalement plat le long de $Y$ aux points de $Z$ (ce qui aura lieu en particulier si $X$ est régulier aux points de $Z$). Si $\sh F$ est normalement plat le long de $Z$ et s'il existe un point $z\in Z$ tel que $\sh F$ soit normalement plat le long de $Y$ au point $z$, alors $\sh F$ est normalement plat le long de $Y$ aux points de $Z$.
\end{enumerate}

Soient $\sh K$ et $\sh L\supset\sh K$ les Idéaux cohérents de $\sh O_X$ définissant respectivement $Y$ et $Z$. Pour tout point $z\in Z$, l'immersion canonique $Z\to Y$ est régulière au point $z$, puisque $Y$ et $Z$ sont réguliers en ce point (19.1.1); autrement dit, il existe une suite $(f_i)_{1\leq i\leq n}$ d'éléments de $\sh O_{X,z}$ qui est $\sh K_z$-régulière et est telle que $\sh L_z=\sh K_z+\sum_i f_i\sh O_{X,z}$.
\begin{enumerate}
  \item[(i)] L'hypothèse signifie que $\operatorname{gr}_{\sh K_z}^*(\sh F_z)$ est un $\sh O_{Y,z}$-module plat; elle entraîne donc que $\operatorname{gr}_{\sh L_z}^*(\sh F_z)$ est un $\sh O_{Z,z}$-module plat, en vertu du fait que dans (19.7.1), $a)$ entraîne $b)$.
  \item[(ii)] L'hypothèse supplémentaire sur $X$ entraîne que la suite $(f_i)$ est $\operatorname{gr}_{\sh K_z}^*(\sh F_z)$-régulière en tout point $z\in Z$ (19.7.2). D'autre part, $Z$ étant régulier et connexe, est intègre, et si $\operatorname{gr}_{\sh K_z}^*(\sh F_z)$ est $\sh O_{Y,z}$-plat en un point $z\in Z$, alors, comme le point générique $\zeta$ de $Z$ est une générisation de $z$, on en conclut que $\operatorname{gr}_{\sh K_\zeta}^*(\sh F_\zeta)$ est $\sh O_{Y,\zeta}$-plat ($0_{\mathrm I}$, 6.3.1). On en déduit que si en tout point $z'$ de $Z$, $\operatorname{gr}_{\sh L_{z'}}^*(\sh F_{z'})$ est un $\sh O_{Z,z'}$-module plat, alors $\operatorname{gr}_{\sh K_{z'}}^*(\sh F_{z'})$ est un $\sh O_{Y,z'}$-module plat, car on peut appliquer l'équivalence de $a)$ et $d)$ dans (19.7.1).
\end{enumerate}
\end{corollary}

\subsection{Propriétés de passage à la limite projective.}
\EGAReaderTarget{ega4.c04.subsec.19-8}{subsection:IV.19.8}
\label{IV.19.8}

Dans ce numéro, les notations et conventions sur les limites projectives sont celles de (8.5.1) et de (8.8.1).

\begin{proposition}[19.8.1]
\label{IV.19.8.1}
Supposons que les morphismes de transition $S_\mu\to S_\lambda$ $(\lambda\leq\mu)$ soient plats, et en outre que l'une des deux hypothèses suivantes soit vérifiée:
\begin{enumerate}
  \item[$1^\circ$] Les préschémas $S_\lambda$ sont localement noethériens.
  \item[$2^\circ$] Les morphismes de transition $S_\mu\to S_\lambda$ sont surjectifs (donc fidèlement plats).
\end{enumerate}

